% Section 2 (was Section 3 before loop0009 item 01; rewritten for readability
% in item 03). Sources: paper2/problem_transformations.md (problems,
% statements, proofs), paper2/equivalences.md (master table, sources, gaps),
% paper2/proof_reductions.md (faults in published proofs). Lean:
% lean/MOSPFormalization/ and Complex/; axioms printed by paper2/axiom_check.lean.
% Lean names in Appendix A; formal definitions of the eight problems of
% Table 2.1, the old 3.3 and the old 3.4 are in Appendix B. The band end
% values (K_2, K_{1,3}, K_{3,3}) are computed by paper2/complex_check.py
% (equivalences.md, "Not formalised, by choice").
\section{The equivalences, proved}
\label{sec:equiv}

This section answers the first half of our question: which rows of the
Linhares--Yanasse table are true, and in what sense. Section~\ref{sec:intro}
separated two readings of ``plus or minus one''. We call a row \emph{exact}
when it meets the fixed-offset reading: its value is pathwidth, or pathwidth
plus one, on every input. We call it a \emph{band} when it meets only the
literal reading: its value stays within a fixed distance of pathwidth, and
both ends of that range occur. We call it \emph{false} when it meets neither.
The answer is that eight rows are exact, two are bands and two are false.
Section~\ref{sec:problems} states the problems. Section~\ref{sec:true} proves
the answer in four steps: the hub theorem, which ties MOSP to pathwidth; the
other exact rows; the bands; and the false rows. Table~\ref{tab:master} and
Figure~\ref{fig:chain} summarise it.

All theorems below are verified in the Lean proof assistant, unless the text
says otherwise; Appendix~\ref{app:lean} lists the formal statements. Where a
result is published, we cite the published proof and do not repeat it.

\paragraph{Notation.} Table~\ref{tab:notation} lists the symbols used in the
rest of the paper. Graphs are finite and simple. A 0/1 matrix $M$ has the
customers as rows and the patterns as columns; in the VLSI problems the rows
are nets and the columns gates. Row $r$ is \emph{active} at position $j$ of a
column order $\pi$ if it has 1s in two columns $x,y$ with
$\pi(x)\le j\le\pi(y)$, possibly $x=y$: its stack is open there. ``$M$ has a
1'' excludes only the all-zero matrix. One letter needs care. In this section
$k$ is the bound in each decision problem, so for pathwidth it bounds a width.
In Section~\ref{sec:search} $k$ bounds the number of open stacks, which is one
more than a width; we say so again there.

\begin{table}[tbp]
\centering\small
\caption{Notation. Each symbol has one meaning throughout the paper; the
symbols of the search are defined in Section~\ref{sec:search-def}.}
\label{tab:notation}
\begin{tabularx}{\linewidth}{@{}>{\raggedright}p{3.2cm}>{\raggedright\arraybackslash}X@{}}
\toprule
Symbol & Meaning \\
\midrule
$G=(V,E)$, $n=|V|$ & a graph, its vertices and edges; in Section~\ref{sec:search} the vertices are the customers, $V=C$ \\
$L$ & a \emph{layout}, a bijection $V\to\{1,\dots,n\}$ \\
$X_1,\dots,X_r$ & the bags of a path decomposition \\
$M$; $C$; $P$ & a 0/1 matrix; its rows (customers, nets); its columns (patterns, gates) \\
$G_M$ & the \emph{MOSP graph}: vertex set $C$, two customers adjacent when some pattern has a 1 in both rows; each pattern becomes a clique \\
$H$ & the \emph{net graph} of a circuit: two nets adjacent when they share a gate \\
$\pi$ & an order of the columns \\
$Z(M)$, $t(M)$ & the MOSP value (fewest open stacks) and the gate matrix layout value (fewest tracks) \\
$\pw$, $\vs$ & pathwidth and vertex separation \\
$\theta$, $\ns$, $\mns$, $\es$, $\pes$, $\nu$, $\sbw$, $\cw$, $\mcw$, $\VSG$, $\pla$ & the values of the other problems (Table~\ref{tab:problems}) \\
$k$ & the bound of a decision problem; in Section~\ref{sec:search}, a bound on open stacks, so width $k-1$ \\
$w$ & a width, as in ``pathwidth $w$'' \\
$N[c]$ & the closed neighbourhood of $c$: $c$ and its neighbours \\
$\partial T$, $b(T)$ & the border of a vertex set $T$, the vertices outside $T$ with a neighbour in $T$; $b(T)=|\partial T|$ \\
\bottomrule
\end{tabularx}
\end{table}

\subsection{The problems}
\label{sec:problems}

Each problem is a decision problem with an instance and an integer $k$; its
value is the least $k$ for which the answer is yes. Four of the twelve carry
the rest of the paper, and we define them in full. Figure~\ref{fig:example}
shows each on one instance.

\begin{description}
\item[Pathwidth $\pw(G)$.] Is there a sequence of bags $X_1,\dots,X_r\subseteq V$
  covering $V$, with every edge inside some bag, $X_i\cap X_\ell\subseteq X_j$
  whenever $i\le j\le\ell$, and $|X_i|\le k+1$? \citep[p.~346]{LY2002ref13}.
  The condition on $X_i\cap X_\ell$ is the \emph{interval property}: the bags
  that contain a vertex are consecutive.
\item[Vertex separation $\vs(G)$.] Is there a layout $L$ such that for every
  $i<n$ at most $k$ vertices $u$ with $L(u)\le i$ have a neighbour $v$ with
  $L(v)>i$? \citep[p.~346]{LY2002ref13}.
\item[MOSP $Z(M)$.] Given a 0/1 matrix $M$ of customers by patterns, is there
  an order of the patterns with at most $k$ customers active at every position?
  \citep[eqs.~(1)--(2)]{LY2002}.
\item[Gate matrix layout $t(M)$.] Given a net--gate matrix, is there an order of
  the gates and an assignment of nets to $k$ tracks such that nets active at a
  common position get different tracks? \citep[p.~18]{LY2002ref6};
  \citep[Problem~1]{LY2002ref8}.
\end{description}

The other eight problems are needed only to read the table.
Table~\ref{tab:problems} gives each in one line; Appendix~\ref{app:defs} gives
the formal definitions, with the sources.

\begin{table}[tbp]
\centering\small
\caption{The other eight problems of the Linhares--Yanasse table look
unrelated (two circuit layouts, two search games and four graph problems), yet
each asks for the least $k$ that is possible. One line each. Formal definitions are in
Appendix~\ref{app:defs}.}
\label{tab:problems}
\begin{tabularx}{\linewidth}{@{}>{\raggedright}p{3.6cm}>{\raggedright\arraybackslash}X@{}}
\toprule
Problem & In one line \\
\midrule
One-dimensional logic & Order the gates of a circuit so that the intervals
  spanned by the nets overlap at most $k$ at any point. Its graph is the
  \emph{net graph} $H$, which joins two nets that share a gate. A boundary
  variant pins two gates to the ends. \\
PLA folding $\pla(M)$ & Gate matrix layout in at most $k$ tracks, with at most
  two nets on a track. With no limit per track (\emph{multiple folding}) it is
  gate matrix layout itself. \\
Interval thickness $\theta(G)$ & Add edges to $G$ to make an interval graph
  whose largest clique has at most $k$ vertices. \\
Node search $\ns(G)$ & Searchers are placed on and removed from vertices. An
  edge is cleared when both its ends hold a searcher, and becomes contaminated
  again if a searcher-free path joins it to a contaminated edge. Clear every
  edge with at most $k$ searchers at once. $\mns(G)$ is the same with no
  recontamination allowed. \\
Edge search $\es(G)$ & Node search with a third move: sliding a searcher along
  an edge clears it. $\pes(G)$ is the same over \emph{progressive} strategies,
  which never recontaminate an edge. \\
Narrowness $\nu(G)$ & Add the vertices one at a time in some order, and drop
  each as soon as all its neighbours have been added. Never hold more than $k$. \\
Split bandwidth $\sbw(G)$ & Repeatedly split a vertex into an edge, sharing its
  neighbours between the two ends, so that the result has a layout in which
  adjacent vertices are at most $k$ apart (its \emph{bandwidth}). \\
Edge separation & The table's reference \citep{LY2002ref14} supports three
  readings: cutwidth $\cw(G)$, the most edges crossing a gap of a layout;
  modified cutwidth $\mcw(G)$, the most edges passing strictly over a vertex;
  and Lengauer's vertex separator game $\VSG(G)$, a pebbling game that counts
  the unpebbled vertices adjacent to pebbled ones. \\
\bottomrule
\end{tabularx}
\end{table}

\subsection{What is true}
\label{sec:true}

Table~\ref{tab:master} gives the answer row by row, and
Figure~\ref{fig:chain} draws it as a graph around pathwidth. The rest of this
subsection proves it.

\begin{table}[tbp]
\centering\small
\caption{The Linhares--Yanasse table, settled. Eight rows are exact, two are
bands and two are false, and every true relation is a relation to pathwidth.
\emph{Relation}: what is true and proved in Lean, on the graph the problem
lives on (the MOSP graph or the net graph $H$ for the matrix problems, the input
graph otherwise). \emph{Published source}: the published work the row rests on; the
Lean proofs are listed in Appendix~\ref{app:lean}. \emph{Status}: \emph{exact}, a fixed
offset; \emph{band}, an inequality whose ends both occur; \emph{false}, the
table's claim fails, with a counterexample family proved in Lean.}
\label{tab:master}
\begin{tabularx}{\linewidth}{@{}>{\raggedright}p{2.7cm}>{\raggedright}p{3.4cm}>{\raggedright\arraybackslash}X l@{}}
\toprule
Problem & Published source & Relation & Status \\
\midrule
MOSP & \citet{Yanasse1997a}; \citet{FellowsLangston1989}; \citet{LY2002} &
  $Z(M)=\pw(G_M)+1$ ($M$ has a 1) & exact \\
Gate matrix layout & \citet{LY2002ref6} & $t(M)=Z(M)=\pw(G_M)+1$ & exact \\
One-dim.\ logic & \citet{LY2002ref7} & tracks $=\theta(H)=\pw(H)+1$; boundary variant $\ge\pw(H)+1$ & exact \\
PLA folding & \citet{LY2002ref6} & $\max(\pw+1,\lceil|N|/2\rceil)\le\pla$, gap unbounded; multiple folding $=\pw+1$ & false; variant exact \\
Interval thickness & \citet{LY2002ref6}; \citet{LY2002ref9} & $\theta=\pw+1$ ($V\ne\emptyset$) & exact \\
Node search & \citet{LY2002ref9,LY2002ref10}; \citet{BienstockSeymour1991} & $\ns=\mns=\vs+1=\pw+1=\theta$ ($E\ne\emptyset$) & exact \\
Edge search & \citet{EllisSudboroughTurner1994}; \citet{LY2002ref10} & $\vs\le\es\le\vs+2$, also for $\pes$ & band \\
Narrowness & \citet{LY2002ref11} & $\nu=\pw+1$ ($V\ne\emptyset$) & exact \\
Split bandwidth & \citet{LY2002ref12} & $\pw\le\sbw\le\pw+1$ & band \\
Pathwidth & \citet{LY2002ref13} & the reference & exact \\
Edge separation$^{a}$ & \citet{LY2002ref14} & $\cw,\mcw$ unbounded against $\pw$ (stars); $\VSG=\max(1,\vs)$ & false; VSG exact \\
Vertex separation & \citet{LY2002ref13} & $\vs=\pw$ & exact \\
\bottomrule
\multicolumn{4}{@{}p{\linewidth}@{}}{\footnotesize $^{a}$Also misattributed:
what \citet{LY2002ref14} proves exactly is the vertex separator game, which is
the table's vertex separation row.}
\end{tabularx}
\end{table}

\begin{figure}[tbp]
\centering
\includegraphics[width=\linewidth]{sec3_fig1_chain}
\caption{Pathwidth is the hub of the table: every exact relation is one or two
steps from it, and the only failures are the two false claims, in dark red. Solid edges are exact relations and dashed edges bands, both
proved in Lean; the dotted edge is LaPaugh's $\mathrm{es}=\mathrm{pes}$, the
one relation not proved, which no row needs. Each false claim is refuted by a
counterexample family proved in Lean. The four matrix problems (top left) live
on the MOSP graph $G_M$ or the net graph $H$. Hypotheses are omitted: $M$ has
a 1, $V\neq\emptyset$, or $G$ has an edge, as in Table~\ref{tab:master}. Two
edges that follow from the others, $\mathrm{mns}=\theta$ and
$\mathrm{ns}=\mathrm{vs}+1$, are not drawn.}
\label{fig:chain}
\end{figure}

\paragraph{The hub.} Two theorems tie the table to pathwidth. The first joins
the two graph rows, vertex separation and pathwidth. The second joins MOSP to
pathwidth, and through it all four matrix problems. The second rests on one
lemma about path decompositions. For a vertex $v$, let $\mathrm{first}(v)$ and
$\mathrm{last}(v)$ be the least and greatest indices of bags containing $v$; by
the interval property, $v\in X_i$ exactly when
$\mathrm{first}(v)\le i\le\mathrm{last}(v)$.

\begin{lemma}[cliques sit in a bag]\label{lem:clique}
In a path decomposition of $G$, every clique of $G$ lies inside one bag~\leanref{clique}.
\end{lemma}
\begin{proof}
The bags containing a vertex $u$ of the clique form the interval
$[\mathrm{first}(u),\mathrm{last}(u)]$. Two vertices of the clique are
adjacent, so their edge lies in a bag and their intervals meet. Pairwise
intersecting intervals of a line have a common point (the Helly property): let
$a$ be the largest left end; every right end is at least every left end, so at
least $a$, and $a$ lies in all of them. So all the intervals share an index
$i$, and the clique lies in $X_i$.
\end{proof}

The first theorem is known. It says the two graph rows of the table are the
same quantity.

\begin{theorem}[vertex separation is pathwidth; known]\label{thm:vs}
For every graph, $\vs(G)=\pw(G)$~\leanref{vs}.
\end{theorem}
This is Theorem~3.1 of \citet{LY2002ref13}. Both directions are constructions.
From a layout of separation $k$, each bag is the set counted at a cut together
with the next vertex. From a decomposition, ordering the vertices by
$\mathrm{first}$ gives a layout of the same width.

The second theorem is the one the rest of the paper uses. In
Figure~\ref{fig:example}(d), the sets of customers active at each step of the
good order are the bags of a path decomposition of the MOSP graph. The theorem
says that this always works, in both directions.

\begin{theorem}[MOSP is pathwidth plus one]\label{thm:mosp}
If $M$ has a 1, then $Z(M)=\pw(G_M)+1$~\leanref{mosp}.
\end{theorem}
\begin{proof}
($\le$) Take a decomposition of $G_M$ of width $k$. The rows of each column $p$
form a clique, so by Lemma~\ref{lem:clique} they lie in a bag $X_{\beta(p)}$.
Order the columns by $\beta(p)$. If row $r$ is active at the position of column
$p$, it has columns $x,y$ with $\beta(x)\le\beta(p)\le\beta(y)$, and
$r\in X_{\beta(x)}\cap X_{\beta(y)}$, so the interval property gives
$r\in X_{\beta(p)}$. At most $k+1$ rows are active there.

($\ge$) Take an order with at most $Z=Z(M)$ active rows at every position. The
active sets, with a singleton bag for each row with no 1, form a path
decomposition: rows sharing a column are both active at its position, and a row
is active on an interval of positions. Its width is $Z-1$, using $Z\ge1$.
\end{proof}
For the table, this settles the MOSP row as exact, with offset one, on the MOSP
graph. Read as a cutting problem, Theorem~\ref{thm:mosp} is Proposition~5 of
\citet{Yanasse1997a}, which turns each pattern into a clique on its piece types.
Read as a layout problem, it is Theorem~7 of \citet{FellowsLangston1989}, with
the column-expansion Lemma~4.1 of \citet{FellowsLangston1987}.
\citet[Prop.~2]{LY2002} identify MOSP with gate matrix layout. The proof above
needs no assumption that the instance is reduced.

The graph matters. The MOSP literature also uses a second graph, with the
patterns as vertices, joined when they share a customer, to split an instance
into independent parts \citep{YanasseSenne2010}. It is natural to expect its
pathwidth to measure the open stacks as well; it does not, in either
direction. One pattern shared by four customers who each have a private
pattern has $Z=4$, while its pattern graph is the star $K_{1,4}$, of
pathwidth~1~\leanref{star}. A single customer with many patterns errs the
other way.

\paragraph{The exact core.} The other exact rows attach to
Theorems~\ref{thm:vs} and~\ref{thm:mosp} by published arguments. We do not
reprove them on paper; each item below cites the published proof, and the Lean
proofs are listed in Appendix~\ref{app:lean}. We call the problems that equal
pathwidth or pathwidth plus one on every input the \emph{exact core}: the eight
exact rows, with multiple folding and the vertex separator game.

\begin{theorem}[the exact core]\label{thm:core}
\leavevmode
\begin{enumerate}
\item Gate matrix layout is MOSP on the same matrix: if $M$ has a 1,
  $t(M)=Z(M)=\pw(G_M)+1$. For a fixed gate order the fewest tracks equals the
  most nets active at one position, by the left-edge algorithm
  \citep[p.~31]{LY2002ref6}~\leanref{core-gml}.
\item Interval thickness: if $V\ne\emptyset$, $\theta(G)=\pw(G)+1$
  \citep[Prop.~3.5]{LY2002ref6}~\leanref{core-it}.
\item One-dimensional logic: if every net lies on a gate, the fewest tracks is
  $\theta(H)$, and if some gate connects a net it is $\pw(H)+1$
  \citep[Thm~3]{LY2002ref7}~\leanref{core-logic}.
\item Narrowness: if $V\ne\emptyset$, $\nu(G)=\pw(G)+1$, and for each order
  $\sigma$, $\nu(\sigma)=\vs(\sigma^{\mathrm{rev}})+1$
  \citep[Prop.~3.1]{LY2002ref11}~\leanref{core-narrow}.
\item Node search: for every graph $\ns(G)=\mns(G)$, and if $E\ne\emptyset$,
  $\ns(G)=\vs(G)+1=\pw(G)+1=\theta(G)$
  \citep[Thm]{LY2002ref9}; \citep[Thms~2.3, 4.1]{LY2002ref10}~\leanref{core-ns}.
\item Multiple folding: the fewest tracks is $t(M)=\pw(G_M)+1$
  \citep[Thm~3.14]{LY2002ref6}~\leanref{core-fold}.
\item Lengauer's vertex separator game: $\VSG(G)=\max(1,\vs(G))$
  \citep[Thm~4]{LY2002ref14}~\leanref{core-vsg}.
\end{enumerate}
\end{theorem}
For the table, this makes the gate matrix layout, one-dimensional logic,
interval thickness, node search and narrowness rows exact, and gives an exact
variant of each false row. A single certified pathwidth therefore answers
every problem in the exact core; Section~\ref{sec:dataset} uses this.

\paragraph{Bands.} Two rows are not fixed offsets of pathwidth, but stay within
a fixed distance of it.

\begin{theorem}[bands]\label{thm:bands}
For every graph, $\pw(G)\le\sbw(G)\le\pw(G)+1$ \citep[Thm~8]{LY2002ref12}, and
$\vs(G)\le\es(G)\le\vs(G)+2$, and the same for $\pes$
\citep[Thm~2.1]{EllisSudboroughTurner1994}~\leanref{bands}.
\end{theorem}
Every end of both bands occurs. $K_2$ has $\sbw=\pw=1$ and $\es=\vs=1$; the star
$K_{1,3}$ has $\pw=1$ and $\sbw=2$, and $\vs=1$ with $\es=2$; and $K_{3,3}$ has
$\vs=3$ with $\es=5$. These values are computed by brute force from each
quantity's definition, not proved in Lean.\footnote{$\sbw(K_{1,3})\ge2$ is
also argued by hand: splits of a tree are trees and keep at least three
leaves.} \citet{LY2002ref12} assumes a connected graph with at least two
vertices; the Lean proof needs neither. LaPaugh's equality $\es=\pes$
\citep{LaPaugh1993} is the one relation we state and do not prove, and no row
needs it.

For the table, this means the following. Relative to the number of open
stacks $Z=\pw+1$, split bandwidth lies in $[Z-1,Z]$ and edge search in
$[Z-1,Z+1]$. Both stay within one of $Z$, so they meet the literal reading of
``plus or minus one''. Since both ends occur, neither is a fixed offset, so
neither meets the fixed-offset reading.

\paragraph{False rows.} Two rows fail even the literal reading: no constant
bounds their distance from pathwidth.

\begin{theorem}[two rows are false]\label{thm:false}
\leavevmode
\begin{enumerate}
\item PLA folding: for a matrix with a 1,
  $\max(\pw(G_M)+1,\lceil|N|/2\rceil)\le\pla(M)$, and the gap to $\pw+1$ is
  unbounded. The $5\times5$ identity needs 3 tracks while $\pw(G_M)+1=1$, and the
  path $P_n$ as a matrix, one gate per edge, has $t=2$ and
  $\pla\ge\lceil n/2\rceil$, so the gap is unbounded on connected instances too
  \citep[Prop.~3.15]{LY2002ref6}~\leanref{false-pla}.
\item Edge separation read as cutwidth: $\pw\le\cw$ always, but
  $\cw(K_{1,n})\ge n/2$ and $\mcw(K_{1,n})\ge n/2-1$ while $\pw(K_{1,n})=1$; so
  neither cutwidth reading is within a constant of pathwidth~\leanref{false-cw}.
\end{enumerate}
\end{theorem}
For the table, these rows are false as stated. Each has an exact variant
(Theorem~\ref{thm:core}): the PLA claim holds for multiple folding, and the
edge-separation reference proves exactly the vertex separator game.

Several published equalities need a nondegeneracy hypothesis, and one published
proof has a gap that we repair; Appendix~\ref{app:edge} gives the details.
Appendix~\ref{app:thirteenth} adds a thirteenth problem, outside the table, to
the exact core.

\paragraph{Summary.} Eight of the twelve rows of the Linhares--Yanasse table
are exact, counting the pathwidth row itself, which is the reference. The other
seven are pathwidth, or pathwidth plus one, exactly: MOSP, gate matrix layout,
one-dimensional logic, interval thickness, node search, narrowness and vertex
separation. So are two variants of the false rows, multiple folding and
Lengauer's vertex separator game. Split bandwidth and edge search are bands of
width one and two, within one of the number of open stacks but not at a fixed
offset. Simple PLA folding and edge separation are false as stated. One
theorem is stated and not proved, LaPaugh's $\es=\pes$, and no row needs it.
One of the table's references, Kashiwabara and Fujisawa (1979), is not held; its
row is proved from \citet{LY2002ref6} and \citet{LY2002ref9}.
